% GENERATED FILE. Edit canonical lessons, then run npm run generate:books.
% canonical-source-hash: fnv1a64:172c1ad4a2207dad
% canonical-lessons: FR-C378-echouer, FR-C378-redoubler, FR-C378-concours, FR-C378-moyenne, FR-C378-rattrapage

\chapter{Examens et résultats}
\label{ch:fr-examens-et-resultats}
\hlchaptermodality{378}

\begin{chapteropening}
\textbf{By the end of this chapter:} \emph{I can talk about passing or failing: failing an exam, repeating a year, a competitive exam, the average mark and a resit.}

Complete the last lesson of chapter 378: I can talk about passing or failing: failing an exam, repeating a year, a competitive exam, the average mark and a resit.
\end{chapteropening}

\section[échouer]{échouer — to fail (an exam); to run aground}

\label{lesson:FR-C378-echouer}



\begin{quote}
Before the new one: say the French for by heart, then the French for the textbook.
\end{quote}

\subsection*{You'll want to know: échouer}
\textbf{échouer} — \textquotedblleft{}to fail (an exam); to run aground\textquotedblright{}.

\begin{quote}
Il a échoué à l'examen du permis de conduire.
\end{quote}

\begin{quote}
\emph{He failed his driving test.}
\end{quote}

\emph{Échouer à un examen}; the noun is \emph{un échec}.

\begin{cousinweb}
\emph{échouer} first meant to run a ship aground. The origin is not settled.
\end{cousinweb}

\begin{grammarlens}[title={What you've built}]
One more word for telling, explaining and describing at B1.
\end{grammarlens}

\subsection*{Your turn}
\begin{itemize}
\raggedright
  \item \emph{Say it:} \emph{échouer}
  \item \emph{Say it:} \emph{échouer}, once more
  \item \emph{Say it:} a sentence of your own with \emph{échouer}
  \item \emph{Recall:} say the French for by heart, then the French for the textbook, then say \emph{échouer} again
\end{itemize}

\begin{tcolorbox}[breakable,colback=teal!4,colframe=teal!35!black,title={Before you move on}]
What does échouer mean? (\textquotedblleft{}To fail (an exam); to run aground\textquotedblright{}.) Say it once more.
\end{tcolorbox}
\section[redoubler]{redoubler — to repeat a year (at school); to increase}

\label{lesson:FR-C378-redoubler}



\begin{quote}
Before the new one: say the French for the textbook, then the French for to fail.
\end{quote}

\subsection*{You'll want to know: redoubler}
\textbf{redoubler} — \textquotedblleft{}to repeat a year (at school); to increase\textquotedblright{}.

\begin{quote}
Elle a redoublé sa première année.
\end{quote}

\begin{quote}
\emph{She repeated her first year.}
\end{quote}

Belgium, Switzerland and Québec use the verb without \emph{re-}. \emph{Redoubler d'efforts} — to try twice as hard.

\begin{cousinweb}
\emph{redoubler} is built from \emph{re-} and \emph{doubler} \textquotedblleft{}to double\textquotedblright{}, inherited from Late Latin \emph{duplāre} (EtymDB), from Latin \emph{duplus} \textquotedblleft{}double\textquotedblright{} (Lewis \& Short).
\end{cousinweb}

\begin{grammarlens}[title={What you've built}]
One more word for telling, explaining and describing at B1.
\end{grammarlens}

\subsection*{Your turn}
\begin{itemize}
\raggedright
  \item \emph{Say it:} \emph{redoubler}
  \item \emph{Say it:} \emph{redoubler}, once more
  \item \emph{Say it:} a sentence of your own with \emph{redoubler}
  \item \emph{Recall:} say the French for the textbook, then the French for to fail, then say \emph{redoubler} again
\end{itemize}

\begin{tcolorbox}[breakable,colback=teal!4,colframe=teal!35!black,title={Before you move on}]
What does redoubler mean? (\textquotedblleft{}To repeat a year (at school); to increase\textquotedblright{}.) Say it once more.
\end{tcolorbox}
\section[le concours]{le concours — the competitive exam; the competition}

\label{lesson:FR-C378-concours}



\begin{quote}
Before the new one: say the French for to fail, then the French for to repeat a year.
\end{quote}

\subsection*{You'll want to know: le concours}
\textbf{le concours} — \textquotedblleft{}the competitive exam; the competition\textquotedblright{}.

\begin{quote}
Il prépare le concours pour devenir professeur.
\end{quote}

\begin{quote}
\emph{He is preparing for the competitive exam to become a teacher.}
\end{quote}

Only a set number pass. The same form in the plural.

\begin{cousinweb}
\emph{concours} comes from Latin \emph{concursus} \textquotedblleft{}a running together\textquotedblright{} (EtymDB; Lewis \& Short), from \emph{concurrere} \textquotedblleft{}to run together\textquotedblright{}, built on \emph{currere} \textquotedblleft{}to run\textquotedblright{}.
\end{cousinweb}

\begin{grammarlens}[title={What you've built}]
One more word for telling, explaining and describing at B1.
\end{grammarlens}

\subsection*{Your turn}
\begin{itemize}
\raggedright
  \item \emph{Say it:} \emph{le concours}
  \item \emph{Say it:} \emph{le concours}, once more
  \item \emph{Say it:} a sentence of your own with \emph{le concours}
  \item \emph{Recall:} say the French for to fail, then the French for to repeat a year, then say \emph{le concours} again
\end{itemize}

\begin{tcolorbox}[breakable,colback=teal!4,colframe=teal!35!black,title={Before you move on}]
What does le concours mean? (\textquotedblleft{}The competitive exam; the competition\textquotedblright{}.) Say it once more.
\end{tcolorbox}
\section[la moyenne]{la moyenne — the average; the pass mark}

\label{lesson:FR-C378-moyenne}



\begin{quote}
Before the new one: say the French for to repeat a year, then the French for the competitive exam.
\end{quote}

\subsection*{You'll want to know: la moyenne}
\textbf{la moyenne} — \textquotedblleft{}the average; the pass mark\textquotedblright{}.

\begin{quote}
J'ai eu la moyenne en maths : dix sur vingt.
\end{quote}

\begin{quote}
\emph{I got the pass mark in maths: ten out of twenty.}
\end{quote}

In France, marks are out of twenty; \emph{avoir la moyenne} is to get ten.

\begin{cousinweb}
\emph{moyenne} is the feminine of \emph{moyen} \textquotedblleft{}middle, average\textquotedblright{}, inherited through Old French \emph{moien} (EtymDB) from Latin \emph{mediānus} \textquotedblleft{}in the middle\textquotedblright{}, from \emph{medius} \textquotedblleft{}middle\textquotedblright{} (Lewis \& Short).
\end{cousinweb}

\begin{grammarlens}[title={What you've built}]
One more word for telling, explaining and describing at B1.
\end{grammarlens}

\subsection*{Your turn}
\begin{itemize}
\raggedright
  \item \emph{Say it:} \emph{la moyenne}
  \item \emph{Say it:} \emph{la moyenne}, once more
  \item \emph{Say it:} a sentence of your own with \emph{la moyenne}
  \item \emph{Recall:} say the French for to repeat a year, then the French for the competitive exam, then say \emph{la moyenne} again
\end{itemize}

\begin{tcolorbox}[breakable,colback=teal!4,colframe=teal!35!black,title={Before you move on}]
What does la moyenne mean? (\textquotedblleft{}The average; the pass mark\textquotedblright{}.) Say it once more.
\end{tcolorbox}
\section[le rattrapage]{le rattrapage — the resit (exam); catching up}

\label{lesson:FR-C378-rattrapage}



\begin{quote}
Before the new one: say the French for the competitive exam, then the French for the average.
\end{quote}

\subsection*{You'll want to know: le rattrapage}
\textbf{le rattrapage} — \textquotedblleft{}the resit (exam); catching up\textquotedblright{}.

\begin{quote}
Il a échoué à l'examen, mais il y a un rattrapage en juillet.
\end{quote}

\begin{quote}
\emph{He failed the exam, but there is a resit in July.}
\end{quote}

Belgium: \emph{la seconde session}. From \emph{rattraper}, to catch up.

\begin{cousinweb}
\emph{rattrapage} is built from \emph{rattraper} \textquotedblleft{}to catch again, to catch up\textquotedblright{} and \emph{-age} (EtymDB). \emph{rattraper} is \emph{re-} + \emph{attraper} \textquotedblleft{}to catch\textquotedblright{}, which rests on \emph{trappe} \textquotedblleft{}a trap\textquotedblright{}, usually traced to a Frankish word.
\end{cousinweb}

\begin{grammarlens}[title={What you've built}]
One more word for telling, explaining and describing at B1.
\end{grammarlens}

\subsection*{Your turn}
\begin{itemize}
\raggedright
  \item \emph{Say it:} \emph{le rattrapage}
  \item \emph{Say it:} \emph{le rattrapage}, once more
  \item \emph{Say it:} a sentence of your own with \emph{le rattrapage}
  \item \emph{Recall:} say the French for the competitive exam, then the French for the average, then say \emph{le rattrapage} again
\end{itemize}

\begin{tcolorbox}[breakable,colback=teal!4,colframe=teal!35!black,title={Before you move on}]
What does le rattrapage mean? (\textquotedblleft{}The resit (exam); catching up\textquotedblright{}.) Say it once more.
\end{tcolorbox}
